(1) 横向受力为张力曲率 \(Ty_{xx}\) 与弹性基础力 \(-Ky\)，因此 \(\mu y_{tt}=Ty_{xx}-Ky\)。用 \(y=Ae^{i(kx-\omega t)}\) 代入，得到 \(\boxed{\omega^2=c^2k^2+\omega_0^2}\)，其中 \(c^2=T/\mu\)、\(\omega_0^2=K/\mu\)，相当于一维Klein–Gordon型波。

(2) 当 \(\omega>\omega_0\)，\(k=\sqrt{\omega^2-\omega_0^2}/c\) 是实数，相速 \(v_p=\omega/k\)、群速 \(v_g=d\omega/dk=c^2k/\omega\)，故 \(v_pv_g=c^2\)。当频率低于截止，\(k=i\sqrt{\omega_0^2-\omega^2}/c\) 为纯虚，得到倏逝场，不应将其赋予实数传播群速。

(3) 把方程乘 \(y_t\) 可得 \(\partial_t[\mu y_t^2/2+Ty_x^2/2+Ky^2/2]+\partial_x(-Ty_xy_t)=0\)。对平面波做时间平均，\(\bar e=(\mu\omega^2+Tk^2+K)|A|^2/4=\mu\omega^2|A|^2/2\)，\(\bar S=Tk\omega|A|^2/2\)，所以 \(\bar S/\bar e=(T/\mu)k/\omega=v_g\)。

(4) 有限弦固定端要求 \(Y(0)=Y(L)=0\)，故 \(k_n=n\pi/L\)、\(\boxed{\omega_n=\sqrt{(T/\mu)(n\pi/L)^2+K/\mu}}\)，\(n=1,2,\dots\)。当 \(K\to0\) 时回到普通弦的等间隔谐振频率；当 \(K\) 很大时各低阶频率聚集于截止附近。
